\documentclass[11pt]{article}

\usepackage[T1]{fontenc}

\usepackage[utf8]{inputenc}

\usepackage[margin=1in]{geometry}

\usepackage{graphicx}

\usepackage{amsmath}

\usepackage{amssymb}

\usepackage{booktabs}

\usepackage{array}

\usepackage{tabularx}

\usepackage{caption}

\usepackage{float}

\usepackage{microtype}

\usepackage{mathptmx}

\usepackage{natbib}

\usepackage[hidelinks]{hyperref}

\captionsetup{font=small,labelfont=bf}

\setlength{\parindent}{1.25em}

\setlength{\parskip}{0pt}

\title{Feasibility assessment for studying e-cigarette use and epilepsy in women under 50 using the provided dataset hints}

\author{}

\date{}

\begin{document}

\maketitle

\begin{center}\small\textit{Research Memo}\end{center}

\begin{abstract}

This memo assesses whether the materials available in this run support the research question: among women younger than 50 years in a public U.S. health dataset, what is the prevalence of e-cigarette use, and is e-cigarette use associated with epilepsy? The available evidence did not include a retrievable dataset with both exposure and outcome variables. The inspected BRFSS 2015 diabetes indicators mirror was a diabetes-feature subset rather than a full BRFSS extract, and it did not contain e-cigarette-use or epilepsy variables. Direct retrieval attempts for candidate raw GitHub CSV mirrors failed during this run, and the resulting log is preserved as \ref{tab:artifact-1}. Because no approved empirical results survived validation, this memo reports feasibility findings only and does not present prevalence or association estimates.

\end{abstract}

\section{Introduction}
The user’s question asks for the prevalence of e-cigarette use among women under age 50 and whether e-cigarette use is associated with epilepsy. Addressing that question requires a dataset containing both a measure of e-cigarette use and an epilepsy variable, along with sufficient demographic information to restrict the sample to women younger than 50. In this run, the available materials pointed to a BRFSS 2015 diabetes indicators subset and to the full BRFSS questionnaire as a contextual reference. The central task was therefore to determine whether the hinted dataset could support the analysis, and if not, whether a retrievable alternative was available. The inspection showed that the accessible BRFSS-derived subset was not suitable for the stated question \cite{ref1}, while the full BRFSS questionnaire confirmed that tobacco-related content exists in the broader survey instrument \cite{ref2}. Because retrieval of a usable raw CSV mirror failed, no original analytic result could be produced.

\section{Methods}
I evaluated the question against the actual materials available in this run. First, I inspected the provided dataset hints and repository metadata to determine whether the referenced BRFSS 2015 mirror was a full survey extract or a smaller derived subset. The repository description indicated a diabetes indicators dataset with features such as BMI, age, income, smoking, physical activity, general health, high blood pressure, and cholesterol, but not e-cigarette use or epilepsy \cite{ref1}. Second, I consulted the BRFSS 2015 questionnaire only to confirm that the full survey includes tobacco-related content, which distinguishes the questionnaire from the smaller diabetes-indicators subset \cite{ref2}. Third, I attempted reproducible retrieval of several candidate raw GitHub URLs for the BRFSS 2015 diabetes indicators CSV and documented the outcomes in a machine-readable log preserved as \ref{tab:artifact-1}. These retrieval attempts failed in this run \cite{ref3}. Because no accessible dataset in the provided materials both contained the relevant variables and was retrievable here, I did not estimate prevalence, fit an association model, or compute any descriptive statistics.

\section{Results}
No approved empirical results were produced in this run. The key empirical pattern from the retrieval log is that all candidate direct-download attempts for raw GitHub CSV mirrors failed, as summarized in \ref{tab:artifact-1}. The metadata review also showed that the accessible BRFSS 2015 diabetes indicators repository is not a suitable source for this question because it is a diabetes-focused feature subset rather than a full BRFSS extract and lacks e-cigarette-use and epilepsy variables \cite{ref1}. As a result, the run did not yield a usable analytic sample, a prevalence estimate for women under 50, or any association result linking e-cigarette use to epilepsy.

\begin{table}[t]
\centering
\caption{Machine-readable log of candidate raw GitHub dataset URLs attempted and the retrieval failure message from this run.}
\label{tab:artifact-1}
\texttt{tables/table\_1.json}
\end{table}

\section{Discussion}
The main conclusion from this run is feasibility-related: the provided materials do not support a valid answer to the stated research question. The dataset hint appears to point to a BRFSS-derived diabetes indicators subset, which is not designed for analyses of e-cigarette use and epilepsy \cite{ref1}. The questionnaire reference confirms that BRFSS as a survey family can include tobacco-related items, but that does not make the specific retrieved subset suitable for the question at hand \cite{ref2}. The failed raw-download attempts further prevented even an exploratory fallback analysis \cite{ref3}. In practical terms, this means the user’s question would need a different data source, ideally a full BRFSS release that includes the relevant modules or another public health survey with both tobacco exposure and neurologic outcomes. The materials reviewed here justify a cautious stop rather than an unsupported estimate.

\section{Limitations}
This run has substantial limitations. First, no valid analytic dataset with both e-cigarette-use and epilepsy variables was available among the materials that could be accessed here. Second, the inspected BRFSS 2015 diabetes mirror was a diabetes-feature subset rather than a full BRFSS extract, and the repository description did not indicate the needed variables \cite{ref1}. Third, direct retrieval of a usable raw CSV mirror failed, so no backup descriptive analysis could be performed \cite{ref3}. Fourth, because there were no approved results, this memo cannot report prevalence, effect estimates, confidence intervals, or p-values. Finally, any answer to the exact question would require switching to a different source that explicitly contains both exposure and outcome variables, such as a full BRFSS release with the relevant modules or another comparable survey.

\bibliographystyle{plainnat}
\bibliography{references}

\end{document}
